\documentclass[aspectratio=169,11pt]{beamer}
% Shared preamble for the talks (Beamer, Metropolis theme)
\usetheme[progressbar=frametitle,numbering=fraction,background=light,block=fill,sectionpage=none]{metropolis}
\usepackage[T1]{fontenc}
\usepackage{FiraSans}
\usepackage[scaled=.85]{FiraMono}
\usepackage{amsmath,amssymb}
\usefonttheme[onlymath]{serif}
\usepackage{booktabs,array,tabularx,multirow}
\usepackage{tikz,pgfplots}
\usepackage{fontawesome5}
\usepackage{qrcode}
\usepackage{xcolor}
\hyphenpenalty=10000\exhyphenpenalty=10000
\usepackage{ragged2e}
\usepackage{adjustbox}
\usepackage{multicol}
\newenvironment{sfig}[1][0.80]{\begin{adjustbox}{max width=\linewidth,max totalheight=#1\textheight,center}}{\end{adjustbox}}
\usetikzlibrary{arrows.meta,positioning,shapes.geometric,shapes.misc,fit,calc,backgrounds,decorations.pathreplacing,patterns,matrix,chains}
\usepgfplotslibrary{fillbetween}
\pgfplotsset{compat=1.17}

% ---- palette (matches vinhqdang.github.io) ----
\definecolor{tdInk}{HTML}{1D2330}
\definecolor{tdMuted}{HTML}{5D6678}
\definecolor{tdLine}{HTML}{D9D5CA}
\definecolor{tdPaper}{HTML}{FBFAF7}
\definecolor{tdBlue}{HTML}{2457D6}
\definecolor{tdGreen}{HTML}{17855A}
\definecolor{tdAmber}{HTML}{B7791F}
\definecolor{tdRed}{HTML}{C8362F}
\definecolor{tdPurple}{HTML}{7A4FD1}
\setbeamercolor{normal text}{fg=tdInk,bg=tdPaper}
\setbeamercolor{alerted text}{fg=tdRed}
\setbeamercolor{example text}{fg=tdGreen}
\setbeamercolor{frametitle}{bg=tdInk,fg=white}
\setbeamercolor{progress bar}{fg=tdBlue,bg=tdLine}
\setbeamercolor{title separator}{fg=tdBlue}
\setbeamercolor{block title}{bg=tdBlue!12,fg=tdInk}
\setbeamercolor{block body}{bg=tdBlue!5,fg=tdInk}
\setbeamercolor{block title alerted}{bg=tdRed!12,fg=tdRed!80!black}
\setbeamercolor{block body alerted}{bg=tdRed!5,fg=tdInk}
\setbeamercolor{block title example}{bg=tdGreen!14,fg=tdGreen!70!black}
\setbeamercolor{block body example}{bg=tdGreen!6,fg=tdInk}
\setbeamerfont{frametitle}{size=\large,series=\bfseries}
\setbeamerfont{footnote}{size=\tiny}
\makeatletter\setlength{\metropolis@frametitle@padding}{1.15ex}\makeatother
\setbeamertemplate{navigation symbols}{}
\setbeamertemplate{footnote}{\hangindent=0.8em\hangafter=1\insertfootnotetext\par}

% ---- bibliography ----
\usepackage[backend=biber,style=authoryear-comp,maxcitenames=2,maxbibnames=4,minbibnames=3,uniquelist=false,giveninits=true,doi=false,url=false,isbn=false,eprint=false,date=year]{biblatex}
\addbibresource{../common/refs.bib}
\addbibresource{../common/own.bib}
\renewcommand*{\nameyeardelim}{\addcomma\space}
\DeclareFieldFormat{doi}{}\renewcommand{\bibfont}{\tiny}
\setlength{\bibitemsep}{1pt}
\newcommand{\src}[1]{\par\vfill{\tiny\color{tdMuted}#1}}   % small source line at slide bottom
\newcommand{\cit}[1]{{\scriptsize\color{tdMuted}\parencite{#1}}}

% ---- tikz styles ----
\tikzset{>={Stealth[length=2.2mm]},
  card/.style={rounded corners=4pt,draw=#1!70!black,fill=#1!10,align=center,inner sep=5pt,font=\small},
  card/.default=tdBlue,
  pill/.style={rounded corners=7pt,draw=none,fill=#1,text=white,inner sep=3pt,font=\scriptsize\bfseries,align=center},
  pill/.default=tdBlue,
  soft/.style={rounded corners=4pt,draw=tdLine,fill=white,align=center,inner sep=4pt,font=\small},
  lbl/.style={font=\scriptsize,text=tdMuted,align=center},
  arr/.style={->,thick,draw=tdMuted},
}
\pgfplotsset{every axis/.append style={tick label style={font=\scriptsize},label style={font=\scriptsize},legend style={font=\scriptsize,draw=none,fill=none}}}

% small helpers
\newcommand{\illus}{{\tiny\color{tdMuted}\itshape illustrative}}
\newcommand{\key}[1]{\textcolor{tdBlue}{\bfseries #1}}
\newcommand{\bad}[1]{\textcolor{tdRed}{\bfseries #1}}
\newcommand{\good}[1]{\textcolor{tdGreen}{\bfseries #1}}
\newcommand{\talkfooter}[1]{\setbeamertemplate{footline}{\hfill\tiny\color{tdMuted}#1\hspace{1.6em}\insertframenumber/\inserttotalframenumber\hspace{0.8em}\vspace{4pt}}}
\newcommand{\fanPts}{(1.57,1.435) (1.05,1.498) (2.55,0.716) (0.82,0.838) (2.21,1.399) (1.70,1.225) (0.77,1.025) (2.12,1.433) (0.71,0.415) (1.90,1.705) (0.81,1.108) (0.87,1.244) (1.87,1.876) (3.08,0.996) (0.97,1.100) (1.27,0.918) (2.48,1.490) (3.44,0.396) (2.33,1.260) (1.79,1.113) (3.53,0.217) (0.74,0.703) (8.98,2.877) (7.27,0.828) (6.83,0.457) (6.75,0.836) (7.33,1.869) (8.85,2.682) (6.94,0.467) (8.14,1.556) (8.32,2.467) (7.52,1.478) (8.04,1.873) (6.59,0.846) (6.58,0.870) (7.02,1.136) (8.44,2.496) (7.68,1.982) (7.34,1.931) (8.16,2.442) (7.76,2.083) (7.30,1.587) (8.78,2.077) (8.50,2.808)}
\newcommand{\fanMean}{(0.00,0.000) (0.20,0.245) (0.40,0.483) (0.60,0.707) (0.80,0.911) (1.00,1.090) (1.20,1.238) (1.40,1.353) (1.60,1.430) (1.80,1.469) (2.00,1.469) (2.20,1.431) (2.40,1.358) (2.60,1.252) (2.80,1.119) (3.00,0.963) (3.20,0.790) (3.40,0.608) (3.60,0.422) (3.80,0.240) (4.00,0.069) (4.20,-0.085) (4.40,-0.216) (4.60,-0.319) (4.80,-0.389) (5.00,-0.423) (5.20,-0.419) (5.40,-0.377) (5.60,-0.296) (5.80,-0.179) (6.00,-0.027) (6.20,0.154) (6.40,0.360) (6.60,0.586) (6.80,0.825) (7.00,1.070) (7.20,1.315) (7.40,1.552) (7.60,1.774) (7.80,1.976) (8.00,2.152) (8.20,2.298) (8.40,2.409) (8.60,2.482) (8.80,2.517) (9.00,2.514) (9.20,2.473) (9.40,2.396) (9.60,2.288) (9.80,2.152) (10.00,1.995)}
\newcommand{\fanUp}{(0.00,0.597) (0.20,0.794) (0.40,0.994) (0.60,1.193) (0.80,1.389) (1.00,1.569) (1.20,1.720) (1.40,1.840) (1.60,1.909) (1.80,1.946) (2.00,1.953) (2.20,1.909) (2.40,1.840) (2.60,1.732) (2.80,1.621) (3.00,1.446) (3.20,1.277) (3.40,1.088) (3.60,0.904) (3.80,0.746) (4.00,0.609) (4.20,0.501) (4.40,0.425) (4.60,0.385) (4.80,0.385) (5.00,0.426) (5.20,0.394) (5.40,0.364) (5.60,0.378) (5.80,0.436) (6.00,0.536) (6.20,0.677) (6.40,0.854) (6.60,1.064) (6.80,1.304) (7.00,1.548) (7.20,1.797) (7.40,2.033) (7.60,2.257) (7.80,2.456) (8.00,2.632) (8.20,2.778) (8.40,2.888) (8.60,2.967) (8.80,2.996) (9.00,2.993) (9.20,2.972) (9.40,2.927) (9.60,2.862) (9.80,2.780) (10.00,2.683)}
\newcommand{\fanLo}{(0.00,-0.597) (0.20,-0.304) (0.40,-0.029) (0.60,0.221) (0.80,0.433) (1.00,0.611) (1.20,0.756) (1.40,0.865) (1.60,0.951) (1.80,0.991) (2.00,0.984) (2.20,0.953) (2.40,0.876) (2.60,0.772) (2.80,0.617) (3.00,0.480) (3.20,0.304) (3.40,0.128) (3.60,-0.060) (3.80,-0.265) (4.00,-0.471) (4.20,-0.671) (4.40,-0.857) (4.60,-1.023) (4.80,-1.163) (5.00,-1.272) (5.20,-1.233) (5.40,-1.118) (5.60,-0.970) (5.80,-0.793) (6.00,-0.591) (6.20,-0.369) (6.40,-0.133) (6.60,0.108) (6.80,0.346) (7.00,0.592) (7.20,0.833) (7.40,1.070) (7.60,1.291) (7.80,1.497) (8.00,1.672) (8.20,1.818) (8.40,1.929) (8.60,1.997) (8.80,2.039) (9.00,2.035) (9.20,1.974) (9.40,1.865) (9.60,1.714) (9.80,1.525) (10.00,1.306)}
\newcommand{\fanUpA}{(0.00,0.462) (0.20,0.707) (0.40,0.945) (0.60,1.169) (0.80,1.373) (1.00,1.552) (1.20,1.700) (1.40,1.815) (1.60,1.892) (1.80,1.931) (2.00,1.931) (2.20,1.893) (2.40,1.820) (2.60,1.714) (2.80,1.581) (3.00,1.425) (3.20,1.252) (3.40,1.070) (3.60,0.884) (3.80,0.702) (4.00,0.531) (4.20,0.377) (4.40,0.246) (4.60,0.143) (4.80,0.073) (5.00,0.039) (5.20,0.043) (5.40,0.085) (5.60,0.166) (5.80,0.283) (6.00,0.435) (6.20,0.616) (6.40,0.822) (6.60,1.048) (6.80,1.287) (7.00,1.532) (7.20,1.777) (7.40,2.014) (7.60,2.236) (7.80,2.438) (8.00,2.614) (8.20,2.760) (8.40,2.871) (8.60,2.944) (8.80,2.979) (9.00,2.976) (9.20,2.935) (9.40,2.858) (9.60,2.750) (9.80,2.614) (10.00,2.457)}
\newcommand{\fanLoA}{(0.00,-0.462) (0.20,-0.217) (0.40,0.021) (0.60,0.245) (0.80,0.449) (1.00,0.628) (1.20,0.776) (1.40,0.891) (1.60,0.968) (1.80,1.007) (2.00,1.007) (2.20,0.969) (2.40,0.896) (2.60,0.790) (2.80,0.657) (3.00,0.501) (3.20,0.328) (3.40,0.146) (3.60,-0.040) (3.80,-0.222) (4.00,-0.393) (4.20,-0.547) (4.40,-0.678) (4.60,-0.781) (4.80,-0.851) (5.00,-0.885) (5.20,-0.881) (5.40,-0.839) (5.60,-0.758) (5.80,-0.641) (6.00,-0.489) (6.20,-0.308) (6.40,-0.102) (6.60,0.124) (6.80,0.363) (7.00,0.608) (7.20,0.853) (7.40,1.090) (7.60,1.312) (7.80,1.514) (8.00,1.690) (8.20,1.836) (8.40,1.947) (8.60,2.020) (8.80,2.055) (9.00,2.052) (9.20,2.011) (9.40,1.934) (9.60,1.826) (9.80,1.690) (10.00,1.533)}
\newcommand{\wealthAlt}{(0,0.000) (1,-0.134) (2,-0.264) (3,-0.157) (4,-0.437) (5,-0.436) (6,-0.710) (7,-0.547) (8,-0.501) (9,-0.305) (10,-0.351) (11,-0.280) (12,-0.297) (13,-0.374) (14,-0.336) (15,-0.480) (16,-0.506) (17,-0.396) (18,-0.369) (19,-0.403) (20,-0.098) (21,-0.071) (22,-0.128) (23,-0.088) (24,-0.136) (25,-0.167) (26,-0.193) (27,0.090) (28,0.113) (29,0.155) (30,0.263) (31,0.545) (32,0.535) (33,0.474) (34,0.816) (35,0.645) (36,0.617) (37,0.723) (38,1.040) (39,0.936) (40,0.639) (41,0.745) (42,0.697) (43,0.666) (44,0.745) (45,0.794) (46,0.851) (47,0.815) (48,1.002) (49,1.218) (50,1.242) (51,1.202) (52,1.317) (53,1.399) (54,1.283) (55,1.243) (56,1.399) (57,1.406) (58,1.383) (59,1.430) (60,1.445) (61,1.464) (62,1.232) (63,1.480) (64,1.525) (65,1.431) (66,1.578) (67,1.638) (68,1.659) (69,1.816) (70,2.130) (71,2.221) (72,2.447) (73,2.748) (74,2.635) (75,2.586) (76,2.707) (77,2.504) (78,2.488) (79,2.676) (80,2.832) (81,2.927) (82,2.699) (83,3.025) (84,3.117) (85,3.236) (86,3.320) (87,3.116) (88,2.949) (89,2.810) (90,3.079) (91,2.990) (92,2.978) (93,2.965) (94,3.059) (95,3.143) (96,3.198) (97,3.053) (98,2.664) (99,2.705) (100,2.759) (101,2.939) (102,2.945) (103,2.885) (104,2.958) (105,2.775) (106,2.659) (107,2.548) (108,2.605) (109,2.897) (110,3.016) (111,3.184) (112,3.234) (113,3.290) (114,3.251) (115,3.234) (116,2.937) (117,2.800) (118,2.652) (119,2.781) (120,2.886) (121,3.055) (122,3.277) (123,3.257) (124,3.414) (125,3.688) (126,3.707) (127,3.868) (128,3.866) (129,3.675) (130,3.710) (131,3.623) (132,3.657) (133,3.616) (134,3.677) (135,3.418) (136,3.709) (137,3.906) (138,3.880) (139,3.781) (140,3.801) (141,3.908) (142,3.983) (143,4.055) (144,3.818) (145,3.690) (146,3.885) (147,3.936) (148,3.757) (149,3.663) (150,3.494) (151,3.381) (152,3.554) (153,3.545) (154,3.712) (155,3.756) (156,3.578) (157,3.512) (158,3.834) (159,3.752) (160,3.633)}
\newcommand{\wealthNullA}{(0,0.000) (1,0.145) (2,0.015) (3,0.024) (4,-0.276) (5,-0.307) (6,-0.142) (7,0.057) (8,-0.111) (9,-0.202) (10,-0.256) (11,-0.629) (12,-0.429) (13,-0.358) (14,-0.308) (15,-0.192) (16,-0.256) (17,-0.246) (18,-0.156) (19,-0.077) (20,-0.160) (21,-0.344) (22,-0.283) (23,-0.250) (24,-0.413) (25,-0.284) (26,-0.220) (27,-0.183) (28,-0.190) (29,-0.136) (30,0.031) (31,-0.030) (32,0.145) (33,0.159) (34,0.196) (35,0.378) (36,0.388) (37,0.427) (38,0.336) (39,0.318) (40,0.320) (41,0.232) (42,0.126) (43,0.229) (44,0.352) (45,0.217) (46,0.115) (47,0.201) (48,0.334) (49,0.184) (50,0.117) (51,0.251) (52,0.489) (53,0.599) (54,0.429) (55,0.581) (56,0.557) (57,0.501) (58,0.535) (59,0.393) (60,0.274) (61,-0.121) (62,0.018) (63,0.033) (64,-0.041) (65,-0.055) (66,-0.080) (67,-0.165) (68,-0.277) (69,-0.337) (70,-0.676) (71,-0.913) (72,-0.974) (73,-0.811) (74,-0.766) (75,-0.680) (76,-0.602) (77,-0.647) (78,-0.623) (79,-0.644) (80,-0.721) (81,-0.743) (82,-0.707) (83,-0.784) (84,-0.898) (85,-0.888) (86,-0.899) (87,-0.889) (88,-0.757) (89,-0.748) (90,-0.765) (91,-0.865) (92,-1.022) (93,-0.967) (94,-0.902) (95,-0.848) (96,-0.794) (97,-0.883) (98,-0.843) (99,-0.833) (100,-0.764) (101,-0.836) (102,-0.698) (103,-0.576) (104,-0.719) (105,-0.748) (106,-0.764) (107,-0.714) (108,-0.563) (109,-0.791) (110,-0.669) (111,-0.597) (112,-0.675) (113,-0.767) (114,-0.704) (115,-0.750) (116,-0.823) (117,-0.886) (118,-0.642) (119,-0.903) (120,-0.676) (121,-0.892) (122,-0.761) (123,-0.838) (124,-0.844) (125,-0.943) (126,-1.004) (127,-0.848) (128,-0.967) (129,-1.277) (130,-1.350) (131,-1.695) (132,-1.813) (133,-1.851) (134,-1.950) (135,-2.038) (136,-1.816) (137,-1.537) (138,-1.340) (139,-1.150) (140,-1.114) (141,-1.232) (142,-1.341) (143,-1.499) (144,-1.621) (145,-1.708) (146,-1.646) (147,-1.485) (148,-1.446) (149,-1.373) (150,-1.478) (151,-1.416) (152,-1.537) (153,-1.563) (154,-1.349) (155,-1.317) (156,-1.362) (157,-1.303) (158,-1.405) (159,-1.408) (160,-1.535)}
\newcommand{\wealthNullB}{(0,0.000) (1,0.082) (2,0.234) (3,0.116) (4,-0.015) (5,0.014) (6,0.221) (7,0.257) (8,0.091) (9,-0.024) (10,-0.047) (11,0.059) (12,0.042) (13,-0.020) (14,-0.064) (15,-0.058) (16,-0.169) (17,-0.184) (18,-0.166) (19,0.028) (20,0.005) (21,0.113) (22,-0.034) (23,0.060) (24,0.066) (25,-0.019) (26,-0.037) (27,0.063) (28,0.006) (29,0.122) (30,0.301) (31,0.493) (32,0.684) (33,0.507) (34,0.535) (35,0.433) (36,0.538) (37,0.533) (38,0.546) (39,0.590) (40,0.499) (41,0.497) (42,0.390) (43,0.641) (44,0.568) (45,0.555) (46,0.314) (47,0.110) (48,-0.012) (49,0.075) (50,0.159) (51,0.180) (52,0.083) (53,0.305) (54,0.159) (55,0.035) (56,0.082) (57,-0.060) (58,-0.041) (59,0.068) (60,0.044) (61,0.082) (62,0.072) (63,-0.029) (64,-0.450) (65,-0.621) (66,-0.637) (67,-0.862) (68,-0.949) (69,-0.937) (70,-0.781) (71,-0.971) (72,-1.054) (73,-1.040) (74,-0.987) (75,-1.179) (76,-0.956) (77,-1.075) (78,-0.951) (79,-0.991) (80,-0.954) (81,-0.755) (82,-0.678) (83,-0.677) (84,-0.821) (85,-0.920) (86,-1.118) (87,-0.873) (88,-1.088) (89,-1.434) (90,-1.238) (91,-1.088) (92,-1.150) (93,-1.057) (94,-1.178) (95,-1.267) (96,-1.217) (97,-1.182) (98,-1.256) (99,-1.062) (100,-0.982) (101,-0.979) (102,-1.146) (103,-1.144) (104,-1.200) (105,-1.225) (106,-1.137) (107,-1.226) (108,-1.207) (109,-1.119) (110,-1.089) (111,-0.807) (112,-1.048) (113,-1.071) (114,-1.240) (115,-1.264) (116,-0.954) (117,-1.181) (118,-1.235) (119,-1.166) (120,-1.169) (121,-1.067) (122,-1.283) (123,-1.173) (124,-1.312) (125,-1.486) (126,-1.428) (127,-1.382) (128,-1.493) (129,-1.614) (130,-1.656) (131,-1.674) (132,-1.712) (133,-1.859) (134,-2.176) (135,-2.022) (136,-2.016) (137,-2.193) (138,-2.413) (139,-2.519) (140,-2.589) (141,-2.882) (142,-2.717) (143,-2.679) (144,-2.738) (145,-2.709) (146,-2.696) (147,-2.576) (148,-2.573) (149,-2.427) (150,-2.243) (151,-2.301) (152,-2.377) (153,-2.508) (154,-2.555) (155,-2.667) (156,-2.683) (157,-2.660) (158,-2.608) (159,-2.775) (160,-2.652)}
\newcommand{\wealthNullC}{(0,0.000) (1,0.050) (2,0.012) (3,0.460) (4,0.510) (5,0.517) (6,0.558) (7,0.469) (8,0.352) (9,0.405) (10,0.325) (11,0.416) (12,0.368) (13,0.105) (14,-0.094) (15,-0.231) (16,-0.522) (17,-0.552) (18,-0.611) (19,-0.503) (20,-0.640) (21,-0.691) (22,-0.632) (23,-0.551) (24,-0.710) (25,-0.521) (26,-0.626) (27,-0.482) (28,-0.543) (29,-0.619) (30,-0.809) (31,-0.846) (32,-0.539) (33,-0.646) (34,-0.455) (35,-0.294) (36,-0.366) (37,-0.233) (38,-0.063) (39,-0.124) (40,-0.216) (41,-0.083) (42,-0.053) (43,0.218) (44,0.258) (45,0.081) (46,0.055) (47,-0.126) (48,-0.042) (49,0.035) (50,-0.040) (51,-0.181) (52,-0.077) (53,-0.148) (54,-0.000) (55,-0.154) (56,-0.210) (57,-0.304) (58,-0.465) (59,-0.112) (60,0.099) (61,-0.006) (62,0.083) (63,-0.182) (64,-0.030) (65,0.100) (66,0.034) (67,-0.033) (68,-0.018) (69,0.026) (70,-0.043) (71,-0.178) (72,-0.232) (73,-0.310) (74,-0.221) (75,-0.228) (76,-0.214) (77,-0.460) (78,-0.531) (79,-0.603) (80,-0.693) (81,-0.727) (82,-0.786) (83,-1.013) (84,-0.902) (85,-0.795) (86,-0.883) (87,-0.815) (88,-0.961) (89,-1.151) (90,-1.197) (91,-1.001) (92,-1.033) (93,-1.259) (94,-1.243) (95,-1.339) (96,-1.454) (97,-1.422) (98,-1.677) (99,-1.938) (100,-2.154) (101,-2.021) (102,-2.023) (103,-2.122) (104,-2.230) (105,-2.155) (106,-1.999) (107,-2.113) (108,-2.091) (109,-1.782) (110,-1.713) (111,-1.771) (112,-1.739) (113,-1.737) (114,-1.659) (115,-1.624) (116,-1.704) (117,-1.734) (118,-1.695) (119,-1.802) (120,-1.805) (121,-1.943) (122,-1.786) (123,-1.585) (124,-1.717) (125,-1.613) (126,-1.603) (127,-1.603) (128,-1.649) (129,-1.625) (130,-1.500) (131,-1.680) (132,-1.553) (133,-1.602) (134,-1.764) (135,-1.783) (136,-1.708) (137,-1.733) (138,-1.705) (139,-1.763) (140,-1.848) (141,-1.724) (142,-1.863) (143,-1.739) (144,-1.705) (145,-1.755) (146,-1.798) (147,-1.850) (148,-1.867) (149,-1.885) (150,-1.916) (151,-2.192) (152,-2.491) (153,-2.457) (154,-2.554) (155,-2.644) (156,-2.679) (157,-2.841) (158,-2.643) (159,-2.655) (160,-2.690)}
\newcommand{\wealthNullD}{(0,0.000) (1,0.263) (2,0.099) (3,0.122) (4,0.069) (5,-0.065) (6,-0.160) (7,-0.247) (8,-0.279) (9,-0.249) (10,-0.136) (11,-0.130) (12,-0.008) (13,0.027) (14,0.016) (15,-0.038) (16,-0.189) (17,-0.420) (18,-0.456) (19,-0.504) (20,-0.489) (21,-0.473) (22,-0.662) (23,-0.781) (24,-0.953) (25,-0.851) (26,-0.893) (27,-0.905) (28,-1.049) (29,-1.241) (30,-1.344) (31,-1.385) (32,-1.354) (33,-1.578) (34,-1.564) (35,-1.703) (36,-1.799) (37,-1.870) (38,-1.954) (39,-1.924) (40,-1.840) (41,-1.825) (42,-1.923) (43,-1.898) (44,-1.878) (45,-1.796) (46,-1.790) (47,-2.027) (48,-2.153) (49,-2.233) (50,-2.143) (51,-2.190) (52,-2.101) (53,-2.093) (54,-2.274) (55,-2.297) (56,-2.254) (57,-2.231) (58,-2.450) (59,-2.424) (60,-2.528) (61,-2.828) (62,-2.930) (63,-3.212) (64,-3.404) (65,-3.248) (66,-2.960) (67,-2.853) (68,-2.826) (69,-2.611) (70,-2.596) (71,-2.750) (72,-3.094) (73,-3.006) (74,-3.325) (75,-3.721) (76,-3.834) (77,-3.854) (78,-3.994) (79,-3.888) (80,-3.928) (81,-3.930) (82,-3.901) (83,-4.005) (84,-4.006) (85,-3.733) (86,-3.882) (87,-3.889) (88,-3.849) (89,-3.947) (90,-3.981) (91,-4.230) (92,-4.314) (93,-4.336) (94,-4.219) (95,-4.124) (96,-4.226) (97,-4.397) (98,-4.187) (99,-4.098) (100,-4.373) (101,-4.444) (102,-4.587) (103,-4.693) (104,-4.942) (105,-5.023) (106,-5.117) (107,-5.174) (108,-5.065) (109,-5.062) (110,-5.246) (111,-5.498) (112,-5.592) (113,-5.469) (114,-5.373) (115,-5.475) (116,-5.349) (117,-5.319) (118,-5.423) (119,-5.392) (120,-5.427) (121,-5.598) (122,-5.679) (123,-5.683) (124,-5.536) (125,-5.696) (126,-5.689) (127,-5.711) (128,-5.742) (129,-5.925) (130,-5.785) (131,-5.640) (132,-5.456) (133,-5.419) (134,-5.250) (135,-5.203) (136,-5.097) (137,-5.164) (138,-5.158) (139,-5.224) (140,-5.385) (141,-5.466) (142,-5.636) (143,-5.805) (144,-5.879) (145,-5.981) (146,-5.734) (147,-5.525) (148,-5.407) (149,-5.327) (150,-5.462) (151,-5.449) (152,-5.660) (153,-5.765) (154,-5.851) (155,-5.922) (156,-5.832) (157,-5.958) (158,-5.665) (159,-5.754) (160,-5.666)}
\newcommand{\wealthNullE}{(0,0.000) (1,-0.078) (2,-0.340) (3,-0.168) (4,0.236) (5,0.283) (6,0.446) (7,0.610) (8,0.710) (9,0.810) (10,0.956) (11,0.901) (12,0.662) (13,0.399) (14,0.474) (15,0.518) (16,0.494) (17,0.148) (18,0.094) (19,0.065) (20,-0.096) (21,-0.224) (22,0.060) (23,0.230) (24,0.323) (25,0.345) (26,0.262) (27,0.191) (28,0.270) (29,0.130) (30,0.078) (31,0.291) (32,0.088) (33,-0.043) (34,0.190) (35,0.224) (36,0.229) (37,0.091) (38,-0.129) (39,0.022) (40,-0.365) (41,-0.319) (42,-0.418) (43,-0.451) (44,-0.767) (45,-0.660) (46,-0.595) (47,-0.573) (48,-0.533) (49,-0.745) (50,-0.740) (51,-0.691) (52,-0.722) (53,-1.190) (54,-1.225) (55,-1.230) (56,-1.409) (57,-1.420) (58,-1.256) (59,-1.288) (60,-1.121) (61,-1.218) (62,-1.342) (63,-1.324) (64,-1.471) (65,-1.732) (66,-1.704) (67,-1.724) (68,-1.660) (69,-1.803) (70,-1.886) (71,-1.836) (72,-1.925) (73,-1.976) (74,-1.963) (75,-2.090) (76,-2.075) (77,-2.305) (78,-2.254) (79,-1.892) (80,-1.821) (81,-1.897) (82,-2.148) (83,-2.333) (84,-2.394) (85,-2.524) (86,-2.473) (87,-2.472) (88,-2.383) (89,-2.434) (90,-2.340) (91,-2.165) (92,-2.313) (93,-2.348) (94,-2.386) (95,-2.419) (96,-2.276) (97,-2.343) (98,-2.490) (99,-2.700) (100,-2.928) (101,-2.962) (102,-2.855) (103,-2.867) (104,-3.000) (105,-3.043) (106,-3.192) (107,-3.067) (108,-3.069) (109,-2.881) (110,-2.806) (111,-2.735) (112,-2.887) (113,-2.834) (114,-2.907) (115,-3.003) (116,-2.995) (117,-2.783) (118,-2.729) (119,-2.461) (120,-2.382) (121,-2.591) (122,-2.682) (123,-2.431) (124,-2.363) (125,-2.366) (126,-2.444) (127,-2.435) (128,-2.389) (129,-2.461) (130,-2.482) (131,-2.456) (132,-2.621) (133,-2.669) (134,-2.784) (135,-2.885) (136,-3.022) (137,-3.122) (138,-3.249) (139,-3.017) (140,-3.052) (141,-3.013) (142,-3.136) (143,-3.128) (144,-3.078) (145,-2.866) (146,-2.635) (147,-2.707) (148,-2.469) (149,-2.468) (150,-2.307) (151,-2.326) (152,-2.234) (153,-2.257) (154,-2.455) (155,-2.384) (156,-2.435) (157,-2.461) (158,-2.434) (159,-2.371) (160,-2.034)}
\newcommand{\wealthNullF}{(0,0.000) (1,-0.233) (2,-0.044) (3,-0.014) (4,-0.109) (5,0.000) (6,-0.153) (7,0.138) (8,0.203) (9,0.180) (10,0.272) (11,0.098) (12,0.158) (13,0.279) (14,0.400) (15,0.457) (16,0.302) (17,0.333) (18,0.138) (19,0.156) (20,0.063) (21,-0.050) (22,0.045) (23,0.303) (24,0.497) (25,0.438) (26,0.374) (27,0.238) (28,0.151) (29,0.160) (30,0.212) (31,0.145) (32,0.137) (33,-0.027) (34,-0.107) (35,0.059) (36,0.230) (37,0.288) (38,0.288) (39,0.307) (40,0.203) (41,0.035) (42,0.092) (43,-0.067) (44,-0.027) (45,-0.196) (46,-0.220) (47,-0.015) (48,0.050) (49,0.022) (50,0.043) (51,-0.177) (52,-0.067) (53,-0.186) (54,-0.309) (55,-0.420) (56,-0.427) (57,-0.535) (58,-0.637) (59,-0.593) (60,-0.586) (61,-0.440) (62,-0.351) (63,-0.209) (64,-0.376) (65,-0.381) (66,-0.329) (67,-0.228) (68,-0.286) (69,-0.221) (70,-0.094) (71,-0.233) (72,-0.363) (73,-0.488) (74,-0.475) (75,-0.638) (76,-0.653) (77,-0.732) (78,-0.880) (79,-1.133) (80,-0.852) (81,-0.703) (82,-0.583) (83,-0.807) (84,-0.709) (85,-0.772) (86,-0.955) (87,-1.047) (88,-1.126) (89,-0.931) (90,-0.929) (91,-0.952) (92,-0.908) (93,-0.862) (94,-1.043) (95,-0.908) (96,-0.896) (97,-0.781) (98,-0.577) (99,-0.637) (100,-0.350) (101,-0.410) (102,-0.464) (103,-0.769) (104,-0.604) (105,-0.647) (106,-0.738) (107,-0.828) (108,-0.771) (109,-0.645) (110,-0.567) (111,-0.618) (112,-0.723) (113,-0.988) (114,-1.090) (115,-1.047) (116,-1.117) (117,-0.880) (118,-0.769) (119,-0.825) (120,-0.775) (121,-0.764) (122,-0.881) (123,-0.940) (124,-1.155) (125,-1.213) (126,-1.093) (127,-1.171) (128,-0.961) (129,-1.003) (130,-0.920) (131,-0.875) (132,-0.976) (133,-1.140) (134,-1.099) (135,-1.011) (136,-1.029) (137,-1.064) (138,-1.450) (139,-1.551) (140,-1.446) (141,-1.324) (142,-1.291) (143,-1.408) (144,-1.309) (145,-1.363) (146,-1.274) (147,-1.101) (148,-1.238) (149,-1.289) (150,-1.270) (151,-1.234) (152,-1.235) (153,-1.137) (154,-1.030) (155,-1.118) (156,-1.147) (157,-1.026) (158,-1.180) (159,-1.269) (160,-1.557)}
\newcommand{\wealthNullG}{(0,0.000) (1,-0.021) (2,-0.128) (3,-0.099) (4,0.051) (5,0.243) (6,0.302) (7,0.286) (8,0.229) (9,0.147) (10,-0.014) (11,-0.246) (12,-0.509) (13,-0.395) (14,-0.426) (15,-0.318) (16,-0.527) (17,-0.516) (18,-0.541) (19,-0.475) (20,-0.529) (21,-0.428) (22,-0.320) (23,-0.167) (24,-0.463) (25,-0.643) (26,-0.479) (27,-0.162) (28,-0.116) (29,-0.179) (30,-0.042) (31,-0.034) (32,-0.026) (33,0.131) (34,0.120) (35,-0.122) (36,-0.231) (37,-0.170) (38,-0.299) (39,-0.269) (40,-0.360) (41,-0.307) (42,-0.394) (43,-0.294) (44,-0.361) (45,-0.522) (46,-0.473) (47,-0.444) (48,-0.445) (49,-0.663) (50,-0.740) (51,-0.632) (52,-0.614) (53,-0.789) (54,-0.643) (55,-0.606) (56,-0.629) (57,-0.707) (58,-0.743) (59,-0.698) (60,-1.021) (61,-1.087) (62,-1.272) (63,-1.235) (64,-1.240) (65,-1.256) (66,-1.222) (67,-1.021) (68,-0.795) (69,-0.979) (70,-0.766) (71,-0.771) (72,-1.015) (73,-0.862) (74,-0.987) (75,-1.091) (76,-0.990) (77,-1.036) (78,-1.034) (79,-0.891) (80,-0.865) (81,-1.230) (82,-1.299) (83,-1.493) (84,-1.678) (85,-1.868) (86,-1.862) (87,-1.868) (88,-1.919) (89,-1.802) (90,-1.805) (91,-1.860) (92,-1.747) (93,-1.694) (94,-1.682) (95,-1.556) (96,-1.472) (97,-1.466) (98,-1.480) (99,-1.557) (100,-1.701) (101,-1.414) (102,-1.389) (103,-1.715) (104,-1.716) (105,-1.833) (106,-1.985) (107,-2.052) (108,-2.034) (109,-1.917) (110,-1.966) (111,-2.011) (112,-1.962) (113,-2.289) (114,-2.519) (115,-2.732) (116,-2.615) (117,-2.683) (118,-2.851) (119,-2.991) (120,-3.185) (121,-3.169) (122,-3.206) (123,-3.165) (124,-3.108) (125,-3.031) (126,-3.007) (127,-3.111) (128,-3.094) (129,-3.129) (130,-3.149) (131,-2.991) (132,-3.117) (133,-3.065) (134,-3.030) (135,-3.046) (136,-3.138) (137,-3.161) (138,-3.216) (139,-3.378) (140,-3.534) (141,-3.777) (142,-3.856) (143,-3.525) (144,-3.684) (145,-3.562) (146,-3.532) (147,-3.392) (148,-3.268) (149,-3.266) (150,-3.209) (151,-3.207) (152,-3.218) (153,-3.244) (154,-3.178) (155,-3.436) (156,-3.491) (157,-3.532) (158,-3.665) (159,-3.693) (160,-3.486)}
\newcommand{\wealthNullH}{(0,0.000) (1,-0.140) (2,-0.322) (3,-0.259) (4,-0.301) (5,-0.098) (6,-0.021) (7,-0.103) (8,-0.014) (9,-0.000) (10,-0.096) (11,0.012) (12,-0.089) (13,0.154) (14,0.261) (15,0.194) (16,0.151) (17,0.251) (18,0.040) (19,0.009) (20,-0.201) (21,-0.201) (22,-0.221) (23,-0.056) (24,-0.105) (25,-0.128) (26,0.031) (27,-0.087) (28,0.054) (29,-0.010) (30,-0.182) (31,-0.096) (32,-0.040) (33,-0.129) (34,0.057) (35,0.242) (36,0.128) (37,0.030) (38,-0.005) (39,-0.051) (40,-0.169) (41,-0.109) (42,-0.129) (43,-0.127) (44,-0.151) (45,-0.148) (46,-0.286) (47,-0.426) (48,-0.506) (49,-0.575) (50,-0.464) (51,-0.435) (52,-0.566) (53,-0.731) (54,-0.577) (55,-0.542) (56,-0.452) (57,-0.628) (58,-0.664) (59,-0.576) (60,-0.402) (61,-0.478) (62,-0.494) (63,-0.537) (64,-0.519) (65,-0.574) (66,-0.540) (67,-0.564) (68,-0.418) (69,-0.478) (70,-0.600) (71,-0.771) (72,-0.755) (73,-1.018) (74,-1.184) (75,-1.157) (76,-0.957) (77,-1.237) (78,-1.349) (79,-1.373) (80,-1.403) (81,-1.205) (82,-1.142) (83,-1.213) (84,-1.060) (85,-0.942) (86,-0.866) (87,-1.008) (88,-0.951) (89,-0.693) (90,-0.707) (91,-0.607) (92,-0.804) (93,-0.892) (94,-0.976) (95,-0.906) (96,-0.694) (97,-0.775) (98,-0.682) (99,-0.614) (100,-0.715) (101,-0.652) (102,-0.675) (103,-0.630) (104,-0.737) (105,-0.662) (106,-0.571) (107,-0.563) (108,-0.637) (109,-0.869) (110,-1.032) (111,-0.903) (112,-0.894) (113,-0.969) (114,-0.854) (115,-0.805) (116,-0.947) (117,-0.995) (118,-0.932) (119,-0.921) (120,-1.009) (121,-1.121) (122,-1.560) (123,-1.679) (124,-1.961) (125,-1.854) (126,-1.748) (127,-1.808) (128,-1.861) (129,-1.983) (130,-1.957) (131,-1.884) (132,-1.810) (133,-1.579) (134,-1.446) (135,-1.383) (136,-1.262) (137,-1.288) (138,-1.333) (139,-1.079) (140,-1.205) (141,-1.362) (142,-1.458) (143,-1.373) (144,-1.536) (145,-1.502) (146,-1.226) (147,-1.197) (148,-1.057) (149,-1.289) (150,-1.373) (151,-1.515) (152,-1.396) (153,-1.364) (154,-1.189) (155,-1.145) (156,-1.062) (157,-1.289) (158,-1.337) (159,-1.076) (160,-1.277)}
\newcommand{\pubsPerYear}{(2016,5) (2017,3) (2018,4) (2019,2) (2020,5) (2021,10) (2022,9) (2023,6) (2024,8) (2025,7) (2026,30)}

% chip node used in the "research map" slides: \chip{col}{y}{colour}{fill\%}{text}; (col*44mm, y)
\newcommand{\chip}[4]{\node[rounded corners=3pt,font=\scriptsize,text width=38mm,align=left,inner sep=3pt,draw=#3!70!black,fill=#3!#4,minimum height=8mm] at (#1*44,#2) }

% compact footnote-style citations at the bottom of a slide: \srcc{free text}{bibkey,bibkey,...}
\DeclareCiteCommand{\scite}
  {\defcounter{maxnames}{3}\defcounter{minnames}{1}}
  {\printnames[given-family]{author} (\printfield{year}). \printfield{title}. %
   \iffieldundef{journaltitle}{\iffieldundef{booktitle}{\iffieldundef{howpublished}{\iffieldundef{institution}{\printfield{note}}{\printfield{institution}}}{\printfield{howpublished}}}{\textit{\printfield{booktitle}}}}{\textit{\printfield{journaltitle}}}.}
  {\ {\color{tdBlue}$\bullet$}\ }
  {}
\setlength{\columnsep}{6pt}
\newcommand{\srcc}[2]{\par\vfill\begingroup\fontsize{5.2}{6.1}\selectfont\color{tdMuted}\raggedright\setlength{\parindent}{0pt}\begin{multicols}{3}\ifx\relax#1\relax\else #1\ \fi\ifx\relax#2\relax\else\scite{#2}\fi\end{multicols}\vspace{-2pt}\endgroup}

\tikzset{chip/.style={rounded corners=3pt,text width=64mm,align=left,inner sep=3pt,minimum height=7.5mm,anchor=north west},
  done/.style={chip,draw=#1!70!black,fill=#1!14},wip/.style={chip,draw=#1!70!black,dashed,fill=#1!3},done/.default=tdBlue,wip/.default=tdBlue}

% research-map list helpers
\newcommand{\pillarhead}[2]{\par\noindent\colorbox{#1}{\textcolor{white}{\strut\ \small\bfseries #2\ }}\par\vspace{1.5mm}}
\newcommand{\pub}[1]{\par\noindent\hangindent=1.1em\hangafter=1 {\color{tdBlue}$\bullet$}\ #1\par\vspace{0.6mm}}
\newcommand{\wip}[1]{\par\noindent\hangindent=1.1em\hangafter=1 {\color{tdMuted}$\circ$}\ #1\par\vspace{0.6mm}}

% five-bar comparison panel: \barpanel{title}{max value}{v1,v2,v3,v4,v5}  (bars: ERM, DP-SGD, PF-LD, DP-EDL, TRACER-Med)
\newcommand{\barpanel}[3]{\begin{tikzpicture}[x=1mm,y=1mm]
  \node[font=\scriptsize\bfseries,anchor=west,text width=60mm] at (0,47) {#1};
  \draw[tdMuted] (0,0) -- (62,0);
  \foreach \bv [count=\bi] in {#3}{%
    \pgfmathsetmacro{\bh}{\bv/#2*36}\pgfmathsetmacro{\bx}{(\bi-1)*12.5+1}%
    \ifnum\bi=5 \def\bcol{tdBlue}\else\def\bcol{tdMuted!55}\fi%
    \fill[\bcol] (\bx,0) rectangle (\bx+9,\bh);%
    \node[font=\tiny,above] at (\bx+4.5,\bh) {\bv};}%
  \foreach \bn [count=\bj] in {ERM,DP-SGD,PF-LD,DP-EDL,TRACER}{\node[font=\tiny,below] at ({(\bj-1)*12.5+5.5},0) {\bn};}
\end{tikzpicture}}

\talkfooter{Trustworthy AI for Smart Systems \,|\, Tay Do University, 15 Dec 2026}

\title{Trustworthy AI for Smart Systems}
\subtitle{From Explainability and Uncertainty to Verification and Resilience}
\author{Dr.\ Quang-Vinh Dang}
\institute{British University Vietnam}
\date{Keynote \,|\, International Conference on Technological Innovation and Sustainable Engineering in the Digital Era\\Tay Do University, 15 December 2026}

\begin{document}
\begin{frame}[plain]\titlepage\end{frame}

%% ===================================================================
%%  A. MOTIVATION  (about 5 min)
%% ===================================================================
\begin{frame}{Smart systems turn data into physical action}
\centering
\begin{sfig}
\begin{tikzpicture}[node distance=8mm]
  \node[card=tdAmber,text width=22mm,minimum height=17mm] (s) {{\Large\faMicrochip}\\\textbf{Sensors}\\[-1pt]{\scriptsize electronics}};
  \node[card=tdBlue,text width=22mm,minimum height=17mm,right=of s] (n) {{\Large\faNetworkWired}\\\textbf{Network, edge}\\[-1pt]{\scriptsize IT}};
  \node[card=tdPurple,text width=22mm,minimum height=17mm,right=of n] (m) {{\Large\faBrain}\\\textbf{AI models}\\[-1pt]{\scriptsize predict, generate}};
  \node[card=tdBlue,text width=22mm,minimum height=17mm,right=of m] (d) {{\Large\faBalanceScale}\\\textbf{Decisions}\\[-1pt]{\scriptsize people, automation}};
  \node[card=tdGreen,text width=22mm,minimum height=17mm,right=of d] (p) {{\Large\faBuilding}\\\textbf{Physical world}\\[-1pt]{\scriptsize construction, plant}};
  \foreach \a/\b in {s/n,n/m,m/d,d/p}{\draw[arr] (\a) -- (\b);}
  \draw[arr,dashed,rounded corners=6pt] (p.south) -- ++(0,-6mm) -| (s.south);
  \node[lbl,fill=tdPaper,inner sep=1pt] at ($(s.south)!0.5!(p.south)+(0,-6mm)$) {actions change what the sensors see next};
  \foreach \x/\t in {s/{drift, faults,\\spoofed signals},n/{intrusions,\\poisoned updates},m/{overconfident or\\fabricated output},d/{unexplained,\\unaccountable},p/{safety and cost\\of errors}}{
    \node[lbl,text=tdRed,below=15mm of \x] {\t};}
\end{tikzpicture}
\end{sfig}

\vspace{3mm}
{\large Every arrow is a place where \bad{trust} can break.}
\end{frame}

\begin{frame}{Accurate is not the same as trustworthy}
\centering
\begin{sfig}
\begin{tikzpicture}[node distance=4mm]
  \node[card=tdAmber,text width=33mm,minimum height=46mm] (a) {{\Large\color{tdAmber}\faBullseye}\\[2pt]{\bfseries\color{tdAmber!70!black}Overconfident}\\[3pt]{\scriptsize Modern networks are often \textbf{poorly calibrated}: they say ``99\%'' and are right far less often.}\\[5pt]{\tiny\color{tdMuted}Guo et al., 2017}};
  \node[card=tdPurple,text width=33mm,minimum height=46mm,right=of a] (b) {{\Large\color{tdPurple}\faCommentDots}\\[2pt]{\bfseries\color{tdPurple!70!black}Fabricating}\\[3pt]{\scriptsize Language models produce fluent, \textbf{unsupported} content: hallucination.}\\[5pt]{\tiny\color{tdMuted}Ji et al., 2023}};
  \node[card=tdRed,text width=33mm,minimum height=46mm,right=of b] (c) {{\Large\color{tdRed}\faBug}\\[2pt]{\bfseries\color{tdRed!70!black}Attackable}\\[3pt]{\scriptsize Tiny input changes, poisoned updates or injected instructions can \textbf{steer} a model.}\\[5pt]{\tiny\color{tdMuted}Szegedy et al., 2014; Greshake et al., 2023}};
  \node[card=tdBlue,text width=33mm,minimum height=46mm,right=of c] (d) {{\Large\color{tdBlue}\faUserSecret}\\[2pt]{\bfseries\color{tdBlue!70!black}Leaky}\\[3pt]{\scriptsize A model can reveal whether a record was in its training data. \textbf{Forgetting} is hard to verify.}\\[5pt]{\tiny\color{tdMuted}Shokri et al., 2017; Bourtoule et al., 2021}};
\end{tikzpicture}
\end{sfig}

\vspace{5mm}
\large Trust needs \key{evidence}, not just an accuracy number.
\srcc{}{guo2017,ji2023,szegedy2014,greshake2023,shokri2017,bourtoule2021}
\end{frame}

\begin{frame}{Four questions a decision-maker should be able to ask}
\centering
\begin{sfig}
\begin{tikzpicture}
  \node[circle,draw=tdInk,fill=tdInk,text=white,font=\bfseries,align=center,inner sep=3pt] (c) {Trustworthy\\smart\\system};
  \node[card=tdBlue,text width=50mm,minimum height=20mm,anchor=east] (e) at (-16mm,19mm) {\textbf{\faSearch\ Explainability}\\{\scriptsize ``\emph{Why} did it decide that?''}\\[-1pt]{\scriptsize attributions $\cdot$ evidence subgraphs}};
  \node[card=tdAmber,text width=50mm,minimum height=20mm,anchor=west] (u) at (16mm,19mm) {\textbf{\faChartArea\ Uncertainty awareness}\\{\scriptsize ``\emph{How sure} are you?''}\\[-1pt]{\scriptsize calibration $\cdot$ conformal $\cdot$ anytime-valid}};
  \node[card=tdGreen,text width=50mm,minimum height=20mm,anchor=east] (v) at (-16mm,-19mm) {\textbf{\faCheckDouble\ Verification}\\{\scriptsize ``Is it \emph{true}?''}\\[-1pt]{\scriptsize grounding $\cdot$ monitors $\cdot$ certified checks}};
  \node[card=tdRed,text width=50mm,minimum height=20mm,anchor=west] (r) at (16mm,-19mm) {\textbf{\faUserShield\ Resilience}\\{\scriptsize ``What if it is \emph{attacked}?''}\\[-1pt]{\scriptsize robust aggregation $\cdot$ privacy $\cdot$ unlearning}};
  \foreach \x in {e,u,v,r}{\draw[tdMuted,thick] (c) -- (\x);}
\end{tikzpicture}
\end{sfig}

\vspace{2mm}
{\small Next: a map of each pillar. Then what we built, and what we learned, in each.}
\end{frame}

\begin{frame}{Policy now asks for the same things}
\begin{columns}[T]
\begin{column}{0.29\textwidth}
\begin{sfig}[0.62]
\begin{tikzpicture}[x=1mm,y=1mm]
  \draw[thick,tdMuted] (3,2) -- (3,-66);
  \foreach \yy/\yr/\tx/\cc in {-4/2019/{EU High-Level Expert Group:\\\emph{Ethics Guidelines for Trustworthy AI}}/tdBlue,
                            -30/2023/{NIST \emph{AI Risk Management}\\\emph{Framework} 1.0}/tdPurple,
                            -56/2024/{EU \emph{AI Act}:\\Regulation (EU) 2024/1689}/tdGreen}{
    \fill[\cc] (3,\yy) circle (1.6mm);
    \node[anchor=west,font=\normalsize\bfseries,text=\cc] at (7,\yy+3) {\yr};
    \node[anchor=north west,font=\scriptsize,align=left,text width=38mm] at (7,\yy) {\tx};}
\end{tikzpicture}
\end{sfig}
\end{column}
\begin{column}{0.69\textwidth}
\footnotesize\renewcommand{\arraystretch}{1.28}
\begin{tabular}{@{}lcccc@{}}
\toprule
\textbf{EU HLEG requirement} & \textbf{Explain} & \textbf{Uncertainty} & \textbf{Verify} & \textbf{Resilience}\\
\midrule
Human agency and oversight & \good{\faCircle} & \good{\faCircle} & \faCircle[regular] & \faCircle[regular]\\
Technical robustness and safety & \faCircle[regular] & \good{\faCircle} & \good{\faCircle} & \good{\faCircle}\\
Privacy and data governance & & & \faCircle[regular] & \good{\faCircle}\\
Transparency & \good{\faCircle} & \faCircle[regular] & \good{\faCircle} & \\
Diversity, non-discrimination, fairness & \good{\faCircle} & \faCircle[regular] & & \faCircle[regular]\\
Societal and environmental well-being & & \faCircle[regular] & & \faCircle[regular]\\
Accountability & \good{\faCircle} & & \good{\faCircle} & \good{\faCircle}\\
\bottomrule
\end{tabular}

{\tiny\color{tdMuted}\faCircle\ main contribution \quad \faCircle[regular]\ partial \quad our own mapping of the seven requirements to the four pillars}
\end{column}
\end{columns}
\srcc{}{hleg2019,nist2023,euaiact2024}
\end{frame}

%% ===================================================================
%%  B. THE BIG PICTURE  (about 8 min)
%% ===================================================================
\begin{frame}{Pillar 1 \,|\, Explainability: show the reasons, and test them}
\begin{columns}[T]
\begin{column}{0.50\textwidth}
\begin{sfig}[0.62]
\begin{tikzpicture}[x=1mm,y=1mm,font=\footnotesize]
  \node[card=tdBlue] (r) at (36,46) {Explanations};
  \node[card=tdGreen,text width=30mm] (a) at (17,27) {\textbf{Built in}\\interpretable models};
  \node[card=tdAmber,text width=30mm] (b) at (57,27) {\textbf{Post hoc}\\explain a black box};
  \draw[tdMuted] (r) -- (a); \draw[tdMuted] (r) -- (b);
  \node[soft,font=\scriptsize,text width=14mm] (a1) at (8,7) {rules,\\trees};
  \node[soft,font=\scriptsize,text width=15mm] (a2) at (27,7) {prototypes, additive};
  \node[soft,font=\scriptsize,text width=15mm] (b1) at (48,7) {local:\\LIME, SHAP};
  \node[soft,font=\scriptsize,text width=15mm] (b2) at (67,7) {global:\\surrogates};
  \foreach \p/\q in {a/a1,a/a2,b/b1,b/b2}{\draw[tdMuted] (\p) -- (\q);}
\end{tikzpicture}
\end{sfig}

\vspace{2mm}
{\small\textbf{Rudin (2019):} for high-stakes decisions prefer models that are interpretable by design. Post hoc explanations can be unfaithful.}
\end{column}
\begin{column}{0.47\textwidth}
\begin{sfig}[0.62]
\begin{tikzpicture}[x=1mm,y=1mm]
  \node[font=\small\bfseries,anchor=west] at (0,48) {Why was this sensor flagged?};
  \node[lbl,anchor=west] at (0,43) {\illus\ contributions to the alert score};
  \draw[tdMuted] (46,2) -- (46,38);
  \foreach \yy/\nm/\vv in {34/{Vibration RMS}/26,26/{Temperature}/18,18/{Humidity}/6,10/{Packet rate}/-9,2/{Load}/-17}{
    \node[font=\scriptsize,anchor=east] at (21,\yy) {\nm};
    \ifdim \vv pt>0pt
      \fill[tdBlue!80] (46,\yy-2.5) rectangle (46+\vv,\yy+2.5);
    \else
      \fill[tdRed!80] (46+\vv,\yy-2.5) rectangle (46,\yy+2.5);
    \fi}
  \node[lbl,text=tdRed,anchor=east] at (46,-6) {toward normal};
  \node[lbl,text=tdBlue,anchor=west] at (46,-6) {toward alarm};
\end{tikzpicture}
\end{sfig}

\vspace{1mm}
{\small An explanation is a \key{claim} about the model. Measure its fidelity.}
\end{column}
\end{columns}
\srcc{Our work:}{lime2016,rudin2019,own2020understandingdecisionmachine,own2026xhbotexplainableheterophily}
\end{frame}

\begin{frame}{Pillar 2 \,|\, Uncertainty: know what you do not know}
\begin{columns}[T]
\begin{column}{0.54\textwidth}
\begin{sfig}[0.66]
\begin{tikzpicture}
\begin{axis}[width=8.6cm,height=5.3cm,xlabel={input $x$},ylabel={prediction},ymin=-3.2,ymax=3.6,xmin=0,xmax=10,xtick=\empty,ytick=\empty,axis lines=left,clip=true]
  \addplot[name path=ue,draw=none] coordinates {\fanUp};
  \addplot[name path=le,draw=none] coordinates {\fanLo};
  \addplot[tdAmber!25] fill between[of=ue and le];
  \addplot[name path=ua,draw=none] coordinates {\fanUpA};
  \addplot[name path=la,draw=none] coordinates {\fanLoA};
  \addplot[tdBlue!30] fill between[of=ua and la];
  \addplot[tdInk,thick] coordinates {\fanMean};
  \addplot[only marks,mark=*,mark size=0.9pt,tdInk!70] coordinates {\fanPts};
  \node[font=\tiny,text=tdBlue!70!black,anchor=west] at (axis cs:0.2,-1.85) {aleatoric: noise in the data};
  \node[font=\tiny,text=tdAmber!60!black,align=center] at (axis cs:5,2.9) {epistemic: no data here};
\end{axis}
\end{tikzpicture}
\end{sfig}
\end{column}
\begin{column}{0.44\textwidth}
\begin{sfig}[0.66]
\begin{tikzpicture}
\begin{axis}[width=5.6cm,height=5.3cm,xlabel={stated confidence},ylabel={actual accuracy},xmin=0,xmax=1,ymin=0,ymax=1,xtick={0,0.5,1},ytick={0,0.5,1},axis lines=left,ybar,bar width=0.06,enlarge x limits=0.05]
  \addplot[fill=tdAmber!70,draw=tdAmber!80!black] coordinates {(0.15,0.12)(0.25,0.20)(0.35,0.28)(0.45,0.37)(0.55,0.44)(0.65,0.50)(0.75,0.57)(0.85,0.63)(0.95,0.70)};
  \addplot[sharp plot,tdInk,dashed,no marks] coordinates {(0,0)(1,1)};
  \node[font=\tiny,anchor=north west,text=tdInk] at (axis cs:0.05,0.98) {\illus\ reliability diagram};
  \node[font=\tiny,anchor=south east,text=tdRed] at (axis cs:0.98,0.05) {gap = overconfidence};
\end{axis}
\end{tikzpicture}
\end{sfig}
\end{column}
\end{columns}
\vspace{-1mm}
{\small \textbf{Two kinds:} irreducible noise (aleatoric) and ignorance that more data can cure (epistemic). Real networks are often \bad{overconfident}, so we must \key{calibrate} and \key{measure} it.}
\srcc{Our work:}{guo2017,kendall2017,ovadia2019,own2026utilizinguncertaintymeasures}
\end{frame}

\begin{frame}{Conformal prediction: a coverage guarantee from one calibration set}
\centering
\begin{sfig}[0.62]
\begin{tikzpicture}[x=1.15mm,y=1mm]
  \draw[->,thick,tdMuted] (0,0) -- (118,0) node[right,font=\small,text=tdInk] {score};
  \foreach \xx in {6,11,14,19,22,27,31,33,38,42,45,49,52,57,61,66,70,75,83,94}{\fill[tdBlue!75] (\xx,0) circle (1.2mm);}
  \node[font=\small,anchor=east] at (-2,0) {calibration};
  \draw[tdRed,very thick] (78,-5) -- (78,7) node[above,font=\small,text=tdRed] {$\hat q$: 95\% of calibration scores lie to its left};
  \node[font=\small,anchor=east] at (-2,-17) {new input};
  \foreach \xx/\nm/\cc in {10/cat/tdGreen,29/dog/tdGreen,100/bird/tdRed}{
     \fill[\cc] (\xx,-17) circle (1.7mm);\node[font=\small,below=1.2mm] at (\xx,-17) {\nm};}
  \draw[dashed,tdRed] (78,-22) -- (78,-12);
  \node[card=tdGreen,font=\normalsize] at (59,-36) {prediction set $\mathcal C(x)=\{\text{cat},\text{dog}\}$};
  \node[card=tdBlue,font=\normalsize,text width=88mm] at (59,-54) {$\mathbb P\big(Y_{\text{new}}\in\mathcal C(X_{\text{new}})\big)\ \ge\ 1-\alpha$\\[1pt]{\small distribution-free, finite-sample, needs only exchangeability}};
\end{tikzpicture}
\end{sfig}

\vspace{2mm}
{\small A model that \key{says ``I am not sure'' with a guarantee}: larger sets when unsure, a single answer when confident.}
\srcc{Our work:}{vovk2005,angelopoulos2023,own2024kernelmethodsconformal,own2026confideconformalfree}
\end{frame}

\begin{frame}{Pillar 3 \,|\, Verification: check the output against evidence}
\begin{sfig}[0.52]
\begin{tikzpicture}[node distance=6mm]
  \node[card=tdPurple,text width=26mm] (o) {\faBrain\\\textbf{AI output}\\{\scriptsize answer, report, plan}};
  \node[card=tdBlue,text width=24mm,right=of o] (c) {\faList\\\textbf{Claims}\\{\scriptsize split into checkable statements}};
  \node[card=tdAmber,text width=34mm,right=of c] (e) {\faDatabase\\\textbf{Evidence}\\{\scriptsize knowledge graph, documents, sensor data}};
  \node[card=tdGreen,text width=26mm,right=of e] (v) {\faBalanceScale\\\textbf{Verdict}\\{\scriptsize supported / contradicted / unverifiable}};
  \foreach \a/\b in {o/c,c/e,e/v}{\draw[arr] (\a) -- (\b);}
  \node[pill=tdGreen,below=6mm of v,xshift=-10mm] (a1) {accept};
  \node[pill=tdAmber,right=2mm of a1] (a2) {correct};
  \node[pill=tdRed,right=2mm of a2] (a3) {abstain};
  \draw[arr] (v) -- (a2.north);
\end{tikzpicture}
\end{sfig}
\vspace{1mm}

\begin{columns}[T]
\column{0.32\textwidth}\centering{\small\textbf{Statistical}}\\[1pt]{\scriptsize runtime monitors, uncertainty gates}
\column{0.32\textwidth}\centering{\small\textbf{Semantic}}\\[1pt]{\scriptsize grounding in retrieved or structured evidence}
\column{0.32\textwidth}\centering{\small\textbf{Formal}}\\[1pt]{\scriptsize SMT checking of small networks}
\end{columns}
\srcc{Our work:}{ji2023,lewis2020,own2026realtimehallucination,own2026hierfinraghierarchicalmultimodal}
\end{frame}

\begin{frame}{Pillar 4 \,|\, Resilience: assume something will go wrong}
\begin{sfig}[0.78]
\begin{tikzpicture}[x=1mm,y=1mm]
  \foreach \i/\t/\c in {0/Data/tdBlue,1/{Training\\(federated)}/tdPurple,2/Model/tdAmber,3/{Deployment\\(agents, tools)}/tdGreen}{
    \node[pill=\c,minimum width=32mm,minimum height=9mm] at (\i*36,50) {\t};}
  \foreach \i/\t in {0/{poisoning,\\backdoors},1/{Byzantine clients,\\gradient leakage},2/{adversarial inputs,\\membership inference},3/{prompt injection,\\tool abuse}}{
    \node[card=tdRed,minimum width=32mm,minimum height=14mm,font=\scriptsize] (a\i) at (\i*36,30) {\t};}
  \foreach \i/\t in {0/{provenance,\\robust cleaning},1/{robust aggregation,\\differential privacy},2/{rejection by uncertainty,\\unlearning audits},3/{runtime monitors,\\attack-path scoring}}{
    \node[card=tdGreen,minimum width=32mm,minimum height=14mm,font=\scriptsize] (d\i) at (\i*36,8) {\t};
    \draw[->,thick,tdMuted] (a\i.south) -- (d\i.north);}
  \node[lbl,anchor=east] at (-18,30) {attacks};
  \node[lbl,anchor=east] at (-18,8) {defences};
\end{tikzpicture}
\end{sfig}

{\small Federated learning keeps data local; it does not make training safe. We still need robust aggregation, privacy accounting and audits.}
\srcc{Our work:}{mcmahan2017,blanchard2017,greshake2023,own2024detectingiotmalware,own2021rightbeforgotten}
\end{frame}

\begin{frame}{Monitoring without fooling yourself: anytime-valid evidence}
\begin{columns}[T]
\begin{column}{0.6\textwidth}
\begin{sfig}[0.7]
\begin{tikzpicture}
\begin{axis}[width=8.8cm,height=5.8cm,xlabel={observations},ylabel={evidence $\log_{10}W_t$},xmin=0,xmax=160,ymin=-1.6,ymax=3.2,axis lines=left,ytick={-1,0,1,2,3},clip=true]
  \foreach \n in {A,B,C,D,E,F,G,H}{\addplot[tdMuted!55,thin,no marks] coordinates {\csname wealthNull\n\endcsname};}
  \addplot[tdRed,very thick,no marks] coordinates {\wealthAlt};
  \draw[tdBlue,dashed,thick] (axis cs:0,1.301) -- (axis cs:160,1.301) node[above left,font=\tiny,text=tdBlue] {alarm threshold $1/\alpha=20$};
  \node[font=\tiny,text=tdRed,anchor=east] at (axis cs:48,2.7) {shifted stream: evidence grows};
  \node[font=\tiny,text=tdMuted,anchor=west,fill=tdPaper,inner sep=1pt] at (axis cs:78,0.55) {healthy streams: rarely cross};
\end{axis}
\end{tikzpicture}
\end{sfig}
\end{column}
\begin{column}{0.38\textwidth}
{\small
Evidence is a \key{betting wealth} $W_t$ that cannot grow in expectation if nothing is wrong.

\vspace{2mm}
\textbf{Ville's inequality:}
\[\mathbb P\big(\exists t:\ W_t\ge 1/\alpha\big)\le\alpha\]
\vspace{-2mm}

\good{Peeking is safe}: look every day, stop whenever you like, the false-alarm budget still holds.}
\vspace{1mm}

{\tiny\illus: simulated Gaussian streams, 8 healthy runs (grey) and one run with a small shift (red).}
\end{column}
\end{columns}
\srcc{}{shafer2021,ramdas2023,vovkwang2021}
\end{frame}

%% ===================================================================
%%  C. MY RESEARCH  (about 14 min)
%% ===================================================================
\begin{frame}{From trust in collaborative systems to trustworthy AI}
\begin{columns}[T]
\begin{column}{0.62\textwidth}
\begin{sfig}[0.7]
\begin{tikzpicture}
\begin{axis}[width=9.4cm,height=5.8cm,ybar,bar width=7pt,ymin=0,ymax=34,xtick={2016,2018,2020,2022,2024,2026},xticklabel style={/pgf/number format/1000 sep=},ylabel={publications per year},enlarge x limits=0.07,axis lines=left,
 xticklabels={2016,2018,2020,2022,2024,2026}]
  \addplot[fill=tdBlue!75,draw=tdBlue!90!black] coordinates {\pubsPerYear};
  \node[font=\tiny,anchor=south] at (axis cs:2026,30.5) {30};
\end{axis}
\end{tikzpicture}
\end{sfig}
{\tiny ORCID record of 89 publications (2027 print dates counted in 2026). Google Scholar: more than 1,000 citations.}
\end{column}
\begin{column}{0.36\textwidth}
\small
\textbf{2014--2018} PhD (Inria, Nancy): \emph{trust assessment} in large-scale collaborative systems\\[2pt]
\textbf{2019} trust scores change cooperation \cit{dang2019toit}\\[2pt]
\textbf{2020--2025} intrusion and malware detection, graphs, explanation, conformal methods\\[2pt]
\textbf{2026} one programme: explainable, uncertainty-aware, verifiable, resilient AI
\end{column}
\end{columns}
\srcc{Selected earlier work:}{own2018trustassessmentlarge,own2019influencetrustscore,own2023detectingintrusionwifi,own2024detectingobfuscatedmalware}
\end{frame}

\begin{frame}{Our work, mapped to the four pillars}
\begin{columns}[T]
\begin{column}{0.49\textwidth}\footnotesize
\pillarhead{tdBlue}{Explainability}
\pub{\textbf{XHBot}: explainable bot detection (EAI AIR 2026)}
\pub{\textbf{Explaining IDS decisions} (2020, 2021)}
\wip{Relation-camouflage theory}
\vspace{0mm}
\pillarhead{tdGreen}{Verification}
\pub{\textbf{CMVKG-Guard}: hallucination correction (Discover AI 2026)}
\pub{HierFinRAG: financial document QA (Informatics 2026)}
\wip{\textbf{CCRC}: certified correction}
\end{column}
\begin{column}{0.49\textwidth}\footnotesize
\pillarhead{tdAmber}{Uncertainty}
\pub{\textbf{Uncertainty-aware IDS} (SN Comput.\ Sci.\ 2026)}
\pub{\textbf{CONFIDE}: conformal causal survival (Mathematics 2026)}
\wip{\textbf{PHASE}: anytime-valid structural monitoring}
\vspace{0mm}
\pillarhead{tdRed}{Resilience}
\pub{\textbf{FORTRESS-FL}: robust federated learning (Array 2026)}
\pub{\textbf{VOUCH}: unlearning certificates (Machine Learning, accepted)}
\pub{\textbf{CAPS}: LLM attack-path scoring (2026)}
\wip{\textbf{SENTRY-Fuse}: agent-trace monitoring}
\end{column}
\end{columns}
\srcc{Filled: published or accepted. Hollow: in progress. Papers shown (selection):}{own2026xhbotexplainableheterophily,own2026confideconformalfree,own2026realtimehallucination,own2026fortressflbyzantine}
\end{frame}

\begin{frame}{Explainability \,|\, XHBot: bots that hide among friends}
\begin{columns}[T]
\begin{column}{0.5\textwidth}
\begin{sfig}[0.62]
\begin{tikzpicture}[x=1mm,y=1mm]
  \foreach \i/\xx/\yy in {1/10/46,2/30/50,3/52/44,4/18/28,5/40/30,6/62/28,7/8/10,8/30/12,9/54/10}{\node[circle,fill=tdBlue!70,minimum size=5mm,inner sep=0] (g\i) at (\xx,\yy) {};}
  \foreach \a/\b in {1/2,2/3,1/4,2/5,3/5,3/6,4/5,5/6,4/7,5/8,6/9,7/8,8/9,5/9}{\draw[tdBlue!50] (g\a) -- (g\b);}
  \foreach \i/\xx/\yy in {1/34/20,2/74/40,3/72/16}{\node[diamond,fill=tdRed,minimum size=6mm,inner sep=0] (b\i) at (\xx,\yy) {};}
  \foreach \a/\b in {1/4,1/5,1/8,1/6,1/9,2/3,2/6,2/5,3/6,3/9,3/5}{\draw[tdRed!70,thick] (b\a) -- (g\b);}
  \node[font=\footnotesize,text=tdBlue!70!black,anchor=west] at (0,58) {\textbf{genuine} users};
  \node[font=\footnotesize,text=tdRed,anchor=west] at (36,58) {\textbf{bots} follow genuine users};
\end{tikzpicture}
\end{sfig}

\vspace{1mm}
{\small\textbf{Relation camouflage:} averaging over neighbours fails when a bot's neighbours look genuine (\emph{heterophily}).}
\end{column}
\begin{column}{0.47\textwidth}
{\small
\textbf{XHBot} couples three ideas:\\[2pt]
\key{SGTR} spectral filter down-weights camouflage edges\\
\key{THCA} separate homophilic, heterophilic and self signals\\
\key{CPD} decouple behaviour from social position\\[3pt]
\textbf{HFE}: instance-level evidence subgraphs and community-level motifs, scored by fidelity and sparsity.}

\vspace{2mm}
\begin{center}
\begin{sfig}[0.62]
\begin{tikzpicture}
  \node[card=tdGreen,minimum width=48mm,inner sep=4pt] {{\huge\bfseries\color{tdGreen!70!black}0.9474}\ {\small F1 on TwiBot-20}\\[1pt]{\scriptsize +9.64\% over recent baselines}};
\end{tikzpicture}
\end{sfig}
\end{center}
\end{column}
\end{columns}
\srcc{Also on TwiBot-22 and Cresci-2017.}{own2026xhbotexplainableheterophily,own2026integratinggraphfeatures,twibot20}
\end{frame}

\begin{frame}{Explainability \,|\, When does looking at neighbours stop helping?}
\begin{columns}[T]
\begin{column}{0.58\textwidth}
\begin{sfig}[0.7]
\begin{tikzpicture}
\begin{axis}[width=8.6cm,height=5.8cm,xlabel={mean neighbourhood size $D$},ylabel={camouflage threshold $\rho^\star$},xmin=2,xmax=100,ymin=0,ymax=0.95,axis lines=left,legend pos=south east,xmode=log,log basis x=10]
  \addplot[domain=2:100,samples=60,thick,tdBlue] {(1-1/sqrt(x))/(1+1)};
  \addplot[domain=2:100,samples=60,thick,tdAmber] {(1-1/sqrt(x))/(1+0.5)};
  \addplot[domain=2:100,samples=60,thick,tdRed] {(1-1/sqrt(x))/(1+0.1)};
  \legend{$c=1$,$c=0.5$,$c=0.1$}
\end{axis}
\end{tikzpicture}
\end{sfig}
\end{column}
\begin{column}{0.4\textwidth}
{\small
Share $\rho$ of a bot's ties aimed at genuine users. Above
\[\rho^\star=\frac{1-D^{-1/2}}{1+c}\]
averaging over neighbours is \bad{worse than ignoring the graph}. $c$ = minority-to-majority ratio.

\vspace{2mm}
\textbf{Prevalence acts as protection:} the more bots there are, the less camouflage each one needs.}
\end{column}
\end{columns}
\srcc{Closed-form threshold from our relation-camouflage manuscript (in progress), not an experiment.}{zhu2020}
\end{frame}

\begin{frame}{Uncertainty \,|\, Intrusion detection: fewer misses, and honest about transfer}
\begin{columns}[T]
\begin{column}{0.52\textwidth}
\begin{sfig}[0.46]
\begin{tikzpicture}[node distance=5mm,font=\footnotesize]
  \node[card=tdBlue,text width=16mm] (f) {network flow};
  \node[card=tdAmber,text width=34mm,right=6mm of f] (u) {classifier score\\+ \textbf{9 uncertainty measures}\\+ distance + Isolation Forest};
  \node[card=tdGreen,text width=22mm,above right=1mm and 6mm of u,yshift=4mm] (a) {confident\\$\rightarrow$ act};
  \node[card=tdRed,text width=22mm,below right=1mm and 6mm of u,yshift=-4mm] (h) {unsure\\$\rightarrow$ analyst};
  \draw[arr] (f) -- (u); \draw[arr] (u) -- (a); \draw[arr] (u) -- (h);
\end{tikzpicture}
\end{sfig}

\vspace{3mm}
{\small \key{Result:} up to \good{50\% fewer false negatives} for minority attack classes on CICIDS 2017, accuracy above 98\%.}
\end{column}
\begin{column}{0.45\textwidth}
\begin{sfig}[0.46]
\begin{tikzpicture}[x=1mm,y=1mm]
  \node[font=\small\bfseries,anchor=west] at (0,60) {Transfer to another network};
  \draw[tdMuted] (6,0) -- (6,52); \draw[tdMuted] (6,0) -- (66,0);
  \foreach \v in {0.5,0.75,1.0}{\pgfmathsetmacro{\yy}{(\v-0.4)*83}\draw[tdLine] (6,\yy) -- (66,\yy);\node[font=\tiny,anchor=east] at (5,\yy) {\v};}
  \node[font=\tiny,rotate=90] at (-3,26) {ROC-AUC};
  \fill[tdBlue!75] (10,0) rectangle (24,{(0.985-0.4)*83}); \node[font=\scriptsize\bfseries,above] at (17,{(0.985-0.4)*83}) {0.985};
  \fill[tdRed!75] (28,0) rectangle (42,{(0.55-0.4)*83}); \node[font=\scriptsize\bfseries,above] at (35,{(0.55-0.4)*83}) {0.54--0.56};
  \fill[tdMuted!60] (46,0) rectangle (60,{(0.5-0.4)*83}); \node[font=\scriptsize\bfseries,above] at (53,{(0.5-0.4)*83}) {0.50};
  \node[font=\scriptsize,align=center] at (17,-6) {same\\dataset}; \node[font=\scriptsize,align=center] at (35,-6) {other\\dataset}; \node[font=\scriptsize,align=center] at (53,-6) {chance};
\end{tikzpicture}
\end{sfig}

\vspace{4mm}
{\small Trained on one network, a detector \bad{falls to chance} on another (7 public datasets).}
\end{column}
\end{columns}
\srcc{Transfer benchmark: our repository cross\_ids.}{own2026utilizinguncertaintymeasures,own2023learningtransferknowledge,cicids2018}
\end{frame}

\begin{frame}{Uncertainty \,|\, Structures: the false-alarm problem of ``monitoring forever''}
\begin{columns}[T]
\begin{column}{0.56\textwidth}
\begin{sfig}[0.7]
\begin{tikzpicture}
\begin{axis}[width=8.4cm,height=5.8cm,xlabel={monitoring epochs (e.g.\ one per few hours)},ylabel={P(at least one false alarm)},xmin=0,xmax=3000,ymin=0,ymax=1.05,axis lines=left]
  \addplot[domain=0:3000,samples=80,very thick,tdRed] {1-(1-0.0027)^x};
  \addplot[domain=0:3000,thick,tdGreen,dashed] {0.05};
  \node[font=\scriptsize,text=tdRed,anchor=east] at (axis cs:2900,0.8) {fixed $3\sigma$ chart};
  \node[font=\scriptsize,text=tdGreen!70!black,anchor=south west] at (axis cs:60,0.07) {anytime-valid budget $\alpha=0.05$};
\end{axis}
\end{tikzpicture}
\end{sfig}
{\tiny\illus: independent epochs, $3\sigma$ limit ($p=0.27\%$ per epoch). Real residuals are correlated, but the message holds.}
\end{column}
\begin{column}{0.42\textwidth}
{\small
Bridges and buildings now produce \key{streams}. A threshold that is fine for one test is applied every few hours for years.

\vspace{2mm}
In our simulations, a $3\sigma$ chart gave its first false alarm after a \bad{median of 2 days}, in \bad{every one} of 300 undamaged records.

\vspace{2mm}
Alarms nobody trusts are eventually ignored, and every false alarm costs an inspection.}
\end{column}
\end{columns}
\srcc{Our PHASE manuscript.}{shafer2021,ramdas2023}
\end{frame}

\begin{frame}{Uncertainty \,|\, PHASE: physics tells the monitor where damage can hide}
\centering
\begin{sfig}[0.22]
\begin{tikzpicture}[node distance=4mm,font=\footnotesize]
  \node[card=tdAmber,text width=26mm] (t) {predict \textbf{temperature, load} effects};
  \node[card=tdBlue,text width=20mm,right=of t] (r) {score the \textbf{residual}};
  \node[card=tdPurple,text width=34mm,right=of r] (c) {keep the \textbf{damage cone}: stiffness loss lowers frequencies};
  \node[card=tdGreen,text width=20mm,right=of c] (e) {betting \textbf{e-process}};
  \node[card=tdRed,text width=14mm,right=of e] (a) {alarm at $1/\alpha$};
  \foreach \p/\q in {t/r,r/c,c/e,e/a}{\draw[arr] (\p) -- (\q);}
\end{tikzpicture}
\end{sfig}
\vspace{2mm}

\begin{columns}[T]
\begin{column}{0.55\textwidth}
\begin{sfig}[0.62]
\begin{tikzpicture}[x=1mm,y=1mm]
  \foreach \yy/\nm/\vv/\lb in {28/{5\% loss, storey 3}/100/{100\% (150/150)},19/{5\% loss after 2 healthy years}/66/{66\%},10/{2\% stiffness loss}/34/{34\%}}{
    \node[font=\scriptsize,anchor=east] at (36,\yy) {\nm};
    \fill[tdGreen!70] (38,\yy-3) rectangle (38+\vv*0.45,\yy+3);
    \node[font=\scriptsize\bfseries,anchor=west] at (38+\vv*0.45+1,\yy) {\lb};}
  \draw[tdMuted] (38,5) -- (38,33);
  \node[lbl] at (60,0) {detection rate in simulated frame records};
\end{tikzpicture}
\end{sfig}
\end{column}
\begin{column}{0.42\textwidth}\raggedright
{\small \good{0 false alarms} in 300 undamaged records at 1 and at 3 years (95\% upper limit 1.2\%). A 5\% loss is found with a median delay of 59 days.}
\end{column}
\end{columns}
\vspace{1mm}

\begin{alertblock}{Where it breaks (we report it)}
{\tiny\raggedright 30--60 days of commissioning destroys false-alarm control. On the real KW51 bridge a reduced version alarmed before the retrofit in most hourly settings (Z24: silent). The guarantee is conditional, and missed detections are not controlled.}
\end{alertblock}
\srcc{PHASE: our manuscript.}{ramdas2023}
\end{frame}

\begin{frame}{Verification \,|\, CMVKG-Guard: correct hallucinations while the model is writing}
\begin{sfig}[0.40]
\begin{tikzpicture}[node distance=5mm]
  \node[card=tdPurple,text width=22mm] (v) {\textbf{VLM} writes token by token};
  \node[card=tdBlue,text width=36mm,right=of v] (k) {\textbf{scene knowledge graph}\\from the image\\+ Wikidata, ConceptNet};
  \node[card=tdAmber,text width=42mm,right=of k] (x) {\textbf{three-layer check}\\visual--text 30\% $\cdot$ graph 40\% $\cdot$ reasoning 30\%};
  \node[card=tdGreen,text width=34mm,right=of x] (s) {\textbf{replace} low-score tokens by grounded ones (look-ahead beam)};
  \foreach \p/\q in {v/k,k/x,x/s}{\draw[arr] (\p) -- (\q);}
\end{tikzpicture}
\end{sfig}
\vspace{2mm}

\begin{columns}[T]
\column{0.3\textwidth}\centering{\Huge\bfseries\color{tdGreen!70!black}84\%}\\[-1pt]{\scriptsize fewer object hallucinations (CHAIR)}
\column{0.3\textwidth}\centering{\Huge\bfseries\color{tdGreen!70!black}91.5\%}\\[-1pt]{\scriptsize F1 on POPE, +5.8 points over the best training-free baseline}
\column{0.3\textwidth}\centering{\Huge\bfseries\color{tdGreen!70!black}12\,ms}\\[-1pt]{\scriptsize median extra latency per token (A100)}
\end{columns}
\vspace{2mm}

{\small \key{Training-free}: it wraps an existing model, so it can be added to a deployed system.}
\srcc{Six benchmarks, eight baselines.}{own2026realtimehallucination,own2025multimodalretrieval,pope2023}
\end{frame}

\begin{frame}{Verification \,|\, CCRC: can we \emph{certify} the correction?}
\begin{columns}[T]
\begin{column}{0.55\textwidth}
\begin{sfig}[0.7]
\begin{tikzpicture}
\begin{axis}[width=8.2cm,height=5.8cm,xlabel={base error rate $\mu$ of the model},ylabel={forced abstention $\ge(\mu-\alpha)/(1-\alpha)$},xmin=0,xmax=0.6,ymin=0,ymax=0.7,axis lines=left]
  \addplot[domain=0.05:0.6,very thick,tdRed,samples=40] {(x-0.05)/(1-0.05)};
  \node[font=\scriptsize,text=tdRed,anchor=east] at (axis cs:0.58,0.28) {accept-or-abstain only};
  \node[font=\scriptsize,anchor=west] at (axis cs:0.02,0.62) {target error $\alpha=0.05$};
\end{axis}
\end{tikzpicture}
\end{sfig}
\end{column}
\begin{column}{0.43\textwidth}
{\small
\textbf{Impossibility:} if a model may only \emph{emit or abstain} its own answer, whenever $\mu>\alpha$ it must refuse at least this fraction, whatever its score.

\vspace{2mm}
\textbf{CCRC:} also allow \emph{substituting} an answer from an \key{independent evidence channel}, and certify the error of everything emitted at level $\alpha$.

\vspace{2mm}
\good{Honest finding:} of a 66.8-point gain, \textbf{64.5 come from the correctness score}, only 2.3 from our procedure.}
\end{column}
\end{columns}
\srcc{CCRC: our work in progress (five benchmarks, 400+ paired splits).}{bates2021,angelopoulos2024crc}
\end{frame}

\begin{frame}{Resilience \,|\, FORTRESS-FL: federated learning that expects bad participants}
\begin{columns}[T]
\begin{column}{0.52\textwidth}
\begin{sfig}[0.62]
\begin{tikzpicture}[x=1mm,y=1mm,font=\footnotesize]
  \foreach \i/\yy/\cc/\tx in {1/48/tdBlue/{operator 1},2/32/tdBlue/{operator 2},3/16/tdBlue/{operator 3},4/0/tdRed/{operator 4 (Byzantine)}}{
    \node[card=\cc,minimum width=30mm] (o\i) at (15,\yy) {\tx};}
  \node[card=tdGreen,minimum width=36mm,minimum height=46mm,text width=34mm,align=center] (s) at (70,24) {\textbf{coordinator}\\[2pt]verify commitments\\spectral clustering\\reputation weights\\DP noise on aggregate\\hash-chained ledger};
  \foreach \i in {1,2,3}{\draw[arr] (o\i.east) -- (s.west |- o\i);}
  \draw[arr,tdRed] (o4.east) -- (s.west |- o4);
  \node[font=\scriptsize,text=tdRed,align=center] at (40,-7) {outlier update is down-weighted};
\end{tikzpicture}
\end{sfig}
\end{column}
\begin{column}{0.45\textwidth}
{\small
Multi-operator networks cannot pool raw data and cannot trust every participant.

\vspace{2mm}
\textbf{Governance layer:} signed forgetting certificates, consent registry, per-round checkpoints.

\vspace{2mm}
\good{Cost:} audited recovery used \textbf{56--97\% less client energy} than retraining for a fixed 60 rounds, but not always less than early-stopped retraining.

\vspace{2mm}
\bad{Limit:} little power for small models; certifies canary memorisation, not parameter erasure.}
\end{column}
\end{columns}
\srcc{Unlearning layer: our follow-up manuscript.}{own2026fortressflbyzantine,own2026stfedxidsspatiotemporal,blanchard2017}
\end{frame}

\begin{frame}{Resilience \,|\, VOUCH: a certificate that data was really forgotten}
\begin{columns}[T]
\begin{column}{0.52\textwidth}
\begin{sfig}[0.62]
\begin{tikzpicture}[x=1mm,y=1mm,font=\footnotesize]
  \node[card=tdBlue,minimum width=18mm] (a) at (6,34) {twin A};
  \node[card=tdBlue,minimum width=18mm] (b) at (6,10) {twin B};
  \node[circle,draw=tdInk,fill=white,font=\scriptsize,align=center,inner sep=1pt] (coin) at (30,22) {fair\\coin};
  \draw[arr] (a) -- (coin); \draw[arr] (b) -- (coin);
  \node[card=tdGreen,minimum width=28mm] (t) at (64,34) {one twin is\\\textbf{trained on}};
  \node[card=tdAmber,minimum width=28mm] (g) at (64,10) {the other is\\a \textbf{ghost}};
  \draw[arr] (coin) -- (t); \draw[arr] (coin) -- (g);
  \node[card=tdPurple,minimum width=40mm] (u) at (46,-14) {after unlearning the twin\\scores must be \textbf{symmetric}};
  \draw[arr,dashed] (t.east) -- ++(7,0) |- (u.east);
\end{tikzpicture}
\end{sfig}
\end{column}
\begin{column}{0.45\textwidth}
{\small
Paired \emph{ghost canaries} give an exact null with no shadow models. A betting e-process then issues a certificate valid at \textbf{any stopping time}, or \bad{revokes} it.

\vspace{2mm}
\textbf{Real models} (GPT-2, Pythia, Phi-1.5; up to 5.1B parameters): un-unlearned models revoked in \good{16/18} runs; retraining, gradient ascent and NPO certified in \good{18/18} at $\varepsilon=0.2$.

\vspace{2mm}
\bad{Scope:} equivalence is supported at $\varepsilon=0.2$, not yet at 0.1 or 0.05.}
\end{column}
\end{columns}
\srcc{Interactive explainer on vinhqdang.github.io.}{dang2026vouch,own2021rightbeforgotten,bourtoule2021}
\end{frame}

\begin{frame}{Resilience \,|\, LLM agents: score the whole attack path, then watch the trace}
\centering
\begin{sfig}[0.62]
\begin{tikzpicture}[node distance=4mm,font=\footnotesize]
  \node[card=tdRed,text width=24mm] (i) {web page with \textbf{injected text}};
  \node[card=tdAmber,text width=20mm,right=of i] (r) {RAG retriever};
  \node[card=tdPurple,text width=18mm,right=of r] (l) {LLM agent};
  \node[card=tdBlue,text width=16mm,right=of l] (q) {SQL tool};
  \node[card=tdGreen,text width=20mm,right=of q] (d) {customer data};
  \foreach \a/\b in {i/r,r/l,l/q,q/d}{\draw[arr,tdRed] (\a) -- (\b);}
\end{tikzpicture}
\end{sfig}
\vspace{1mm}

\begin{columns}[T]
\begin{column}{0.5\textwidth}\raggedright
{\small \textbf{CAPS} multiplies exposure along the path, attenuates it by guardrails, and ranks mitigations by risk removed per unit cost.

\vspace{3mm}
\textbf{Autonomous-agent benchmark:}\\component-wise (CVSS-style) score \bad{85.0}\\CAPS critical-path score \good{51.3}\\[2pt]{\scriptsize chained low-risk steps are neither independent nor additive}}
\end{column}
\begin{column}{0.47\textwidth}
\begin{sfig}[0.62]
\begin{tikzpicture}[x=1mm,y=1mm]
  \node[font=\small\bfseries,anchor=west] at (0,52) {SENTRY-Fuse: recall at 5\% false alarms};
  \draw[tdMuted] (6,0) -- (6,44); \draw[tdMuted] (6,0) -- (66,0);
  \foreach \v in {0,0.5,1.0}{\pgfmathsetmacro{\yy}{\v*42}\draw[tdLine] (6,\yy) -- (66,\yy);\node[font=\tiny,anchor=east] at (5,\yy) {\v};}
  \fill[tdAmber!70] (12,0) rectangle (28,{0.709*42}); \draw[thick,tdAmber!60!black] (20,{(0.709-0.357)*42}) -- (20,{(0.709+0.357)*42}); \node[font=\scriptsize\bfseries,anchor=west] at (29,{0.709*42}) {0.71 $\pm$ 0.36};
  \fill[tdGreen!70] (12,-0.1) rectangle (28,0); 
  \fill[tdGreen!70] (40,0) rectangle (56,{0.898*42}); \draw[thick,tdGreen!60!black] (48,{(0.898-0.02)*42}) -- (48,{(0.898+0.02)*42}); \node[font=\scriptsize\bfseries,anchor=south] at (48,{(0.898+0.02)*42+1}) {0.90 $\pm$ 0.02};
  \node[font=\scriptsize,align=center] at (20,-6) {one signal};\node[font=\scriptsize,align=center] at (48,-6) {SENTRY-Fuse};
\end{tikzpicture}
\end{sfig}

{\scriptsize 90 compromised agent traces (AgentDojo), 10 splits. The gain is mostly \key{stability}.}
\end{column}
\end{columns}
\srcc{SENTRY-Fuse: our work in progress.}{own2026capscompositionalattack,greshake2023,debenedetti2024}
\end{frame}

\begin{frame}{What did \emph{not} work is part of the evidence}
\centering
\begin{sfig}[0.80]
\begin{tikzpicture}[x=1mm,y=1mm]
  \foreach \i/\c/\t/\s in {
    0/tdRed/{Intrusion detectors do not transfer}/{ROC-AUC 0.985 on the same dataset, 0.54--0.56 on another},
    1/tdRed/{A fancier statistic lost to a threshold}/{SENTRY: the e-process was beaten by a fixed threshold, so we dropped it},
    2/tdAmber/{A guarantee that is vacuous when data are scarce}/{private conformal sets are empty below about 1,300 calibration records},
    3/tdAmber/{Short commissioning breaks the guarantee}/{PHASE with 30--60 days of baseline: false-alarm control is lost},
    4/tdBlue/{Writing the proof found our bug}/{SEAL: a composition bound we had violated in code, fixed by proving it},
    5/tdBlue/{Complexity that did not pay}/{TRACER-Med: subgroup-calibrated clipping gave no consistent benefit}}{
    \pgfmathsetmacro{\col}{mod(\i,2)}\pgfmathsetmacro{\row}{floor(\i/2)}
    \node[card=\c,text width=62mm,minimum height=16mm,align=left,font=\footnotesize,anchor=north west] at (\col*70,-\row*21) {\textbf{\t}\\\s};}
\end{tikzpicture}
\end{sfig}

\vspace{2mm}
\large Report them. \key{Trust is a property of the evidence}, and honest limits are evidence.
\end{frame}

%% ===================================================================
%%  D. SYNTHESIS  (about 3 min)
%% ===================================================================
\begin{frame}{A practical checklist for trustworthy smart systems}
\footnotesize\renewcommand{\arraystretch}{1.12}
\begin{tabularx}{\textwidth}{@{}>{\raggedright\bfseries}p{23mm}>{\raggedright}p{36mm}>{\raggedright}X>{\raggedright\arraybackslash}X@{}}
\toprule
\normalsize Pillar & \normalsize Ask the vendor or the team & \normalsize Evidence to demand & \normalsize Starting tools\\
\midrule
\textcolor{tdBlue}{Explain} & Why this decision? & fidelity and sparsity of explanations; interpretable baseline & attributions, evidence subgraphs\\
\textcolor{tdAmber}{Uncertainty} & How sure, and how do you know? & calibration curves; coverage on held-out sites & conformal sets, anytime-valid alarms\\
\textcolor{tdGreen}{Verify} & Is the output supported? & grounding rate; what is checked, what is not & knowledge-graph checks, runtime monitors\\
\textcolor{tdRed}{Resilience} & What if a participant or input is hostile? & red-team results; privacy budget; deletion audit & robust aggregation, DP, unlearning certificates\\
\midrule
\textbf{All four} & Does it still work \emph{here}? & transfer tests on \emph{your} data (not the vendor's) & held-out-site validation\\
\bottomrule
\end{tabularx}
\end{frame}

\begin{frame}{Working together}
\centering
\begin{sfig}[0.7]
\begin{tikzpicture}
  \node[card=tdBlue,text width=40mm,minimum height=38mm] (a) {{\Large\faFlask}\\[2pt]\textbf{Joint research}\\[3pt]{\scriptsize datasets from construction monitoring and IoT; joint papers; shared benchmarks}};
  \node[card=tdGreen,text width=40mm,minimum height=38mm,right=6mm of a] (b) {{\Large\faChalkboardTeacher}\\[2pt]\textbf{Teaching and mentoring}\\[3pt]{\scriptsize faculty seminars; student projects; co-supervision; academic exchange}};
  \node[card=tdAmber,text width=40mm,minimum height=38mm,right=6mm of b] (c) {{\Large\faIndustry}\\[2pt]\textbf{Applied pilots}\\[3pt]{\scriptsize a trustworthy-AI review of a smart-building, plant or IoT deployment}};
\end{tikzpicture}
\end{sfig}
\vspace{3mm}

{\small The most useful contribution you can make: a \key{real stream of data and a real decision} to test against.}
\end{frame}

\begin{frame}{Take-aways}
\Large
\begin{enumerate}
\item \key{Accurate is not trustworthy.} Ask for explanation, uncertainty, verification and resilience.
\item \key{Guarantees must hold while you watch}: calibrate, use conformal and anytime-valid methods.
\item \key{Test where it will run}: transfer, small calibration sets and hostile participants are where claims fail.
\item \key{Publish the failures} too.
\end{enumerate}
\end{frame}

\begin{frame}[plain]
\centering\vfill
{\Huge\bfseries Thank you}\\[6mm]
{\large Questions and discussion}\\[8mm]
\begin{columns}[c]
\column{0.38\textwidth}\raggedleft
\qrcode[height=26mm]{https://vinhqdang.github.io}
\column{0.58\textwidth}\raggedright\small
\textbf{Dr.\ Quang-Vinh Dang}\\British University Vietnam\\[3pt]
\faGlobe\ vinhqdang.github.io\\
\faGithub\ github.com/vinhqdang\\[3pt]
{\scriptsize Interactive teaching demos and the VOUCH explainer are on the website.}
\end{columns}
\vfill
\end{frame}

%% ===================================================================
%%  REFERENCES
%% ===================================================================
\end{document}
